%%%%%%%%%%%%%%%%%%%%%%%%%%%%%%%%%%%%%%%%%%%%%%%%%%%%%%%%%%%%%%%%%%%%%
%% sec_theory_EM_v2.tex
%%
%% Sections 2.2 -- 2.4 for paperH_HLmain.tex
%% RMG Periodic RT-TDDFT -- Paper 2
%% REVISED April 2026: clean separation of EM-field perturbations
%% (gauge-dependent, Track A) from scalar-potential perturbations
%% (gauge-independent, Track B).
%%
%% Subsection structure:
%%   2.2  Interaction with electromagnetic radiation
%%        2.2.1  Plane wave and minimal coupling
%%        2.2.2  Multipole expansion and long-wavelength approximation
%%        2.2.3  Crystal momentum conservation: vertical transitions
%%        2.2.4  Beyond the optical limit: scalar-potential perturbations
%%   2.3  Gauge transformation: length to velocity gauge
%%        2.3.1  Why the length gauge fails for periodic systems
%%        2.3.2  The Goppert-Mayer transformation
%%        2.3.3  The velocity-gauge TDKS Hamiltonian
%%        2.3.4  The perturbative momentum kick
%%   2.4  Vector potential and magnetic perturbations
%%        2.4.1  Electric perturbation: spatially uniform A(t)
%%        2.4.2  Magnetic perturbation: spatially varying A(r)
%%        2.4.3  Scope of the present implementation
%%
%% CONVENTIONS (all from notation.tex):
%%   Atomic units: hbar = m_e = e = 4*pi*eps_0 = 1
%%   Electron charge = -1  (e > 0 is elementary charge magnitude)
%%   Imaginary unit: \imath  (dotless i, consistent with paper 1)
%%   \myblue{} for all annotations and placeholders -- NEVER \myred{}
%%
%% SOURCES:
%%   NEAT-notes-RMG_main.tex lines 116-716 (primary)
%%   NEAT-project-notes_main.tex lines 256-716 (supplementary)
%%   Physics discussion notes, April 2026
%%%%%%%%%%%%%%%%%%%%%%%%%%%%%%%%%%%%%%%%%%%%%%%%%%%%%%%%%%%%%%%%%%%%%

%--------------------------------------------------------------------
\subsection{REMAINING STUFF}
\label{sec:other}
%--------------------------------------------------------------------



To describe the interaction of electrons with an external electromagnetic field, we start from the time-dependent Kohn--Sham (TDKS) equation for the occupied Kohn--Sham orbitals,
\begin{equation}
i\frac{\partial}{\partial t}\psi_{n}(\mathbf r,t) = \hat H_{\mathrm{KS}}(t)\psi_{n}(\mathbf r,t),
\label{eq:tdks_general}
\end{equation}
where, in atomic units, the TDKS Hamiltonian in the presence of external scalar and vector potentials is written as
\begin{equation}
\hat H_{\mathrm{KS}}(t) = \frac{1}{2} \left[ \hat{\mathbf p} + \frac{1}{c}\mathbf A_{\mathrm{ext}}(\mathbf r,t) \right]^2
                       - \Phi_{\mathrm{ext}}(\mathbf r,t) + \hat V_{\mathrm{ion}} + \hat V_{\mathrm{H}}[n](\mathbf r,t) + \hat V_{\mathrm{xc}}[n](\mathbf r,t).
\label{eq:tdks_minimal_coupling}
\end{equation}
Here \(\hat{\mathbf p}=-i\nabla\), \(\Phi_{\mathrm{ext}}(\mathbf r,t)\) and \(\mathbf A_{\mathrm{ext}}(\mathbf r,t)\) are the external scalar and vector potentials, and \(\hat V_{\mathrm{ion}}\), \(\hat V_{\mathrm{H}}[n]\), and \(\hat V_{\mathrm{xc}}[n]\) denote the ionic, Hartree, and exchange-correlation contributions, respectively. In practical implementations based on pseudopotentials, \(\hat V_{\mathrm{ion}}\) may contain both local and nonlocal parts. The corresponding electric and magnetic fields are obtained from
\begin{equation}
\mathbf E(\mathbf r,t)
=
-\nabla \Phi_{\mathrm{ext}}(\mathbf r,t)
-
\frac{1}{c}\frac{\partial \mathbf A_{\mathrm{ext}}(\mathbf r,t)}{\partial t},
\qquad
\mathbf B(\mathbf r,t)
=
\nabla \times \mathbf A_{\mathrm{ext}}(\mathbf r,t).
\label{eq:em_fields_from_potentials}
\end{equation}

For finite systems, such as atoms and molecules, the interaction with light is commonly treated within the long-wavelength, or electric-dipole, approximation. In this limit the spatial phase factor of the electromagnetic wave,
\begin{equation}
e^{i(\mathbf q\cdot \mathbf r-\omega t)},
\end{equation}
varies negligibly over the extent of the system, so that \(\mathbf q\cdot \mathbf r \ll 1\) and the field may be regarded as spatially uniform. The interaction is then conveniently written in the length gauge as
\begin{equation}
\hat V_{\mathrm{ext}}^{(L)}(t)
=
-\mathbf r \cdot \mathbf E(t).
\label{eq:length_gauge_perturbation}
\end{equation}
Equivalently, one may regard this perturbation as arising from the scalar potential
\begin{equation}
\Phi_{\mathrm{ext}}^{(L)}(\mathbf r,t)
=
\mathbf r\cdot \mathbf E(t),
\qquad
\mathbf A_{\mathrm{ext}}^{(L)}(\mathbf r,t)=0.
\label{eq:length_gauge_potentials}
\end{equation}
Within this approximation, the magnetic component of the radiation field is neglected, and the physically relevant quantity is the field polarization direction. Since the wavevector \(\mathbf q\) is no longer retained explicitly, the terms transverse and longitudinal are not the most useful descriptors at this stage.

The length-gauge form of Eq.~\eqref{eq:length_gauge_perturbation} is well suited for finite systems, but becomes problematic for periodic solids. In particular, the scalar potential \(-\mathbf E(t)\cdot\mathbf r\) is not lattice periodic and is therefore incompatible with Bloch-periodic boundary conditions. In real-space representations this issue is often manifested as a discontinuous sawtooth-like potential across the cell boundary. More fundamentally, the response of a periodic solid to a homogeneous electric field is more naturally expressed in terms of macroscopic current or polarization rather than a direct scalar-potential coupling through the position operator.

For this reason, periodic RT-TDDFT is commonly formulated in the velocity gauge. For a spatially homogeneous electric field, a gauge transformation removes the scalar potential and introduces a purely time-dependent vector potential,
\begin{equation}
\Phi_{\mathrm{ext}}^{(V)}(\mathbf r,t)=0,
\qquad
\mathbf A_{\mathrm{ext}}^{(V)}(\mathbf r,t)=\mathbf A(t),
\qquad
\mathbf E(t)
=
-\frac{1}{c}\frac{d\mathbf A(t)}{dt}.
\label{eq:velocity_gauge_uniform}
\end{equation}
The TDKS Hamiltonian then takes the form
\begin{equation}
\hat H_{\mathrm{KS}}^{(V)}(t)
=
\frac{1}{2}
\left[
\hat{\mathbf p}
+
\frac{1}{c}\mathbf A(t)
\right]^2
+
\hat V_{\mathrm{ion}}
+
\hat V_{\mathrm{H}}[n](\mathbf r,t)
+
\hat V_{\mathrm{xc}}[n](\mathbf r,t).
\label{eq:tdks_velocity_gauge}
\end{equation}
This formulation is gauge-equivalent to the dipole length gauge of Eq.~\eqref{eq:length_gauge_perturbation}, but is considerably better suited for periodic systems because the perturbation preserves translational symmetry. It is important to emphasize that the homogeneous-field velocity gauge of Eq.~\eqref{eq:velocity_gauge_uniform} is still a dipole-limit description. Since \(\mathbf A(t)\) carries no spatial dependence, it does not represent a full propagating electromagnetic wave and does not generate a magnetic field, \(\nabla\times\mathbf A(t)=0\).

A more general description of radiative coupling is obtained by retaining the spatial dependence of the external field. In Coulomb gauge,
\begin{equation}
\nabla \cdot \mathbf A_{\mathrm{ext}}(\mathbf r,t)=0,
\label{eq:coulomb_gauge_condition2}
\end{equation}
a transverse electromagnetic plane wave may be written as
\begin{equation}
\mathbf A_{\mathrm{ext}}(\mathbf r,t)
=
\mathbf A_{0}
e^{i(\mathbf q\cdot\mathbf r-\omega t)}
+\mathrm{c.c.},
\qquad
\mathbf q\cdot\mathbf A_{0}=0.
\label{eq:transverse_plane_wave}
\end{equation}
The condition \(\mathbf q\cdot\mathbf A_{0}=0\) expresses the transverse character of the field, namely that the polarization direction is perpendicular to the propagation direction \(\mathbf q\). In this case both electric and magnetic components are present, with
\begin{equation}
\mathbf E(\mathbf r,t)
=
-\frac{1}{c}\frac{\partial \mathbf A_{\mathrm{ext}}(\mathbf r,t)}{\partial t},
\qquad
\mathbf B(\mathbf r,t)
=
\nabla\times \mathbf A_{\mathrm{ext}}(\mathbf r,t).
\label{eq:transverse_wave_fields}
\end{equation}
This finite-\(\mathbf q\) formulation is more general than the dipole approximation because it explicitly retains momentum transfer and spatial phase information. In the long-wavelength limit, \(e^{i\mathbf q\cdot\mathbf r}\approx 1\), it reduces to the homogeneous-field description discussed above.

It is useful to distinguish the transverse plane wave of Eq.~\eqref{eq:transverse_plane_wave} from a scalar-potential perturbation of the form
\begin{equation}
\Phi_{\mathrm{ext}}(\mathbf r,t)
=
\Phi_{0}\,e^{i(\mathbf q\cdot\mathbf r-\omega t)},
\qquad
\mathbf A_{\mathrm{ext}}(\mathbf r,t)=0.
\label{eq:scalar_probe}
\end{equation}
In this case the electric field is
\begin{equation}
\mathbf E(\mathbf r,t)
=
-\nabla \Phi_{\mathrm{ext}}(\mathbf r,t)
=
-i\mathbf q \,\Phi_{0}\,e^{i(\mathbf q\cdot\mathbf r-\omega t)},
\label{eq:longitudinal_field}
\end{equation}
which is parallel to \(\mathbf q\). Such a perturbation is therefore longitudinal and probes the density response of the system at finite momentum transfer. This type of field is natural for the calculation of density-density response functions, loss spectra, and plasmonic excitations, whereas the transverse vector-potential form of Eq.~\eqref{eq:transverse_plane_wave} describes radiative coupling to an external electromagnetic wave.

The distinction between these external-field representations is also reflected in the analysis of RT-TDDFT simulations. For finite systems treated in the dipole approximation, the response is commonly characterized through the time-dependent dipole moment and its Fourier transform, which yield optical absorption spectra and frequency-dependent polarizabilities. In periodic systems formulated in the velocity gauge, the natural observable is instead the time-dependent current density, from which optical conductivity, dielectric response, and related quantities may be extracted. In the exact limit, length-gauge and velocity-gauge descriptions are equivalent, provided that observables and operators are transformed consistently. In practice, however, the current-based formulation is the natural choice for periodic systems under homogeneous fields, while the more general finite-\(\mathbf q\) framework provides a clear route toward momentum-resolved response beyond the dipole limit.

In the present work we employ the homogeneous-field velocity-gauge formulation of Eqs.~\eqref{eq:velocity_gauge_uniform} and \eqref{eq:tdks_velocity_gauge} as the working framework for periodic RT-TDDFT. This choice preserves periodic boundary conditions, provides direct access to current-based observables, and establishes a natural starting point for future extensions to finite-\(\mathbf q\) perturbations and momentum-resolved spectroscopies.



